% Generated from paper/rh/main.tex by scripts/sync_paper.py.
% Edit the modular sources and regenerate this copy.
\documentclass[11pt]{article}

\usepackage[T1]{fontenc}
\usepackage[utf8]{inputenc}
\usepackage[english]{babel}
\usepackage[margin=1in]{geometry}
\usepackage{lmodern}
\usepackage{microtype}
\usepackage{mathtools}
\usepackage{amsmath,amssymb,amsthm}
\usepackage{array}
\usepackage{booktabs}
\usepackage{tabularx}
\usepackage{float}
\usepackage{graphicx}
\usepackage{pdflscape}
\usepackage{caption}
\usepackage{enumitem}
\usepackage{fancyhdr}
\usepackage[numbers]{natbib}
\usepackage{doi}
\usepackage{xurl}
\usepackage{longtable}
\usepackage{needspace}
\usepackage{aliascnt}
\usepackage{hyperref}
\usepackage{cleveref}
\usepackage{orcidlink}
\usepackage{titling}



\DeclareMathOperator{\Fix}{Fix}
\DeclareMathOperator{\Ran}{Ran}
\DeclareMathOperator{\dist}{dist}
\DeclareMathOperator{\trace}{tr}
\DeclareMathOperator{\Real}{Re}

\newcommand{\JL}{J_{L}}
\newcommand{\psimRH}{\psi_{-}}
\newcommand{\Lzeta}{\Lambda_{\zeta}}
\newcommand{\Znt}{Z_{\zeta}^{\mathrm{nt}}}
\newcommand{\AZ}{A_{Z}(\zeta)}
\newcommand{\Sel}{\mathrm{Sel}^{!}_{\zeta,\mathrm{tr}}}
\newcommand{\Itr}{I_{\mathrm{tr}}}
\newcommand{\GamSDTCSelberg}{\Gamma_{\mathrm{SDTC\text{-}Selberg}}}

\newcommand{\Pminus}{P_{-}}
\newcommand{\psim}{\psi_{-}}
\newcommand{\AX}{A_{X}}
\newcommand{\Kminus}{K^{-}}
\newcommand{\Jiso}{J_{\mathrm{iso}}}
\newcommand{\IOL}{\hat{\mathfrak{X}}}
\newcommand{\GamSDTC}{\Gamma_{\mathrm{SDTC}}}

\newcommand{\SiblingDC}{[1]}

\newtheorem{theorem}{Theorem}[section]
\newtheorem*{theorem*}{Theorem}
\newcommand{\sharedtheorem}[2]{\newaliascnt{#1}{theorem}\newtheorem{#1}[#1]{#2}\aliascntresetthe{#1}}
\sharedtheorem{lemma}{Lemma}
\sharedtheorem{proposition}{Proposition}
\sharedtheorem{corollary}{Corollary}
\theoremstyle{definition}
\sharedtheorem{definition}{Definition}
\sharedtheorem{obligation}{Hypothesis}
\sharedtheorem{classification}{Record}
\sharedtheorem{remark}{Remark}
\newtheorem*{example}{Example}
\floatstyle{ruled}
\newfloat{algorithm}{tbp}{loa}
\floatname{algorithm}{Algorithm}

\urlstyle{same}
\DeclareRobustCommand{\leanid}[1]{{\def\UrlBreaks{\do\.\do\_}\def\UrlBigBreaks{}\def\UrlNoBreaks{}\nolinkurl{#1}}}
\makeatletter
\g@addto@macro\UrlBreaks{\do\-\do\:}
\makeatother
\captionsetup{
  font=small,
  labelfont=bf,
  labelsep=period
}
\captionsetup[table]{skip=6pt,position=top}
\captionsetup[figure]{skip=6pt}
\captionsetup[algorithm]{skip=6pt,position=top}
\crefname{algorithm}{algorithm}{algorithms}
\Crefname{algorithm}{Algorithm}{Algorithms}
\crefname{theorem}{theorem}{theorems}
\Crefname{theorem}{Theorem}{Theorems}
\crefname{lemma}{lemma}{lemmas}
\Crefname{lemma}{Lemma}{Lemmas}
\crefname{proposition}{proposition}{propositions}
\Crefname{proposition}{Proposition}{Propositions}
\crefname{corollary}{corollary}{corollaries}
\Crefname{corollary}{Corollary}{Corollaries}
\crefname{definition}{definition}{definitions}
\Crefname{definition}{Definition}{Definitions}
\crefname{remark}{remark}{remarks}
\crefname{appendix}{Appendix}{Appendices}
\Crefname{appendix}{Appendix}{Appendices}
\crefname{obligation}{hypothesis}{hypotheses}
\Crefname{obligation}{Hypothesis}{Hypotheses}
\crefname{classification}{record}{records}
\Crefname{classification}{Record}{Records}
\Crefname{remark}{Remark}{Remarks}
\setlist[enumerate]{leftmargin=*, itemsep=2pt, topsep=4pt}
\setlength{\droptitle}{-3.2em}
\pretitle{\begin{center}\LARGE\bfseries}
\posttitle{\par\end{center}\vspace{0.5em}}
\preauthor{\begin{center}\normalsize}
\postauthor{\par\end{center}\vspace{-0.4em}}
\predate{\begin{center}\small}
\postdate{\par\end{center}\vspace{0.6em}}
\setlength{\emergencystretch}{2em}
\clubpenalty=10000
\widowpenalty=10000
\displaywidowpenalty=10000
\interfootnotelinepenalty=10000
\allowdisplaybreaks[2]
\fancypagestyle{firstpage}{\fancyhf{}
  \fancyfoot[C]{\scriptsize Automorph Inc., Wilmington, DE, USA \quad \texttt{ioannis@automorph.io} \quad DOI: 10.5281/zenodo.23086975\\[1pt]\textcopyright\ 2026 \quad CC-BY 4.0 \quad Preprint --- not peer reviewed.}
  \renewcommand{\headrulewidth}{0pt}
  \renewcommand{\footrulewidth}{0pt}
}

\title{Riemann Hypothesis via Self-Dual Trace Confinement: A Conditional Closure}
\author{Ioannis Tsiokos\,\orcidlink{0009-0009-7659-5964}}
\date{Version 3: October 1, 2026\\[2pt] \footnotesize (Version 1: June 17, 2026)}
\hypersetup{
pdftitle={Riemann Hypothesis via Self-Dual Trace Confinement: A Conditional Closure},
  pdfauthor={Ioannis Tsiokos},
  hidelinks
}

\begin{document}

\maketitle
\thispagestyle{firstpage}

\begin{abstract}
For any positive integer weights \(m_\rho\), the Riemann hypothesis is
equivalent to the vanishing of \(\AZ=\sum_\rho m_\rho(\Real\rho-\frac12)^2\)
over the nontrivial zeros.  A companion paper proves a
duality-confinement theorem: for a strongly measurable observable with
values in a real Hilbert space that intertwines an involution with a
linear isometric involution and has square-integrable odd part, the
Gram operator of the odd part vanishes whenever it is dominated by
positive trace-class operators of vanishing trace.  This paper applies
the theorem to the reflection \(s\mapsto1-\bar s\) in the critical line
and measures what the application is worth.  A named hypothesis
\(\GamSDTCSelberg\) on a shell of represented zeros, supplying
domination records and a bridge to the zeros' coordinates, implies that
every represented zero has coordinate \(\frac12\); together with a map
from coordinates to the reals that sends \(\frac12\) to \(\frac12\) and
carries the represented zeros onto the real parts of all nontrivial
zeros, it implies the Riemann hypothesis.  Conversely, the hypothesis
holds exactly when every represented zero has coordinate \(\frac12\), so
on the bare shell of all nontrivial zeros with their real parts it is
equivalent to the Riemann hypothesis; so is the concrete version on
finite windows of actual zeros that are closed under
\(s\mapsto1-\bar s\), carry unit weights and eventually contain every
zero.  The conditional theorem is therefore a reformulation, not
evidence.  By elementary linear algebra we also show that, for windows
closed under \(s\mapsto1-s\) with symmetric positive weights that
eventually contain every zero, several natural bounds for the window
sums are again equivalent to the Riemann hypothesis, and that on a
nonempty window operator-level transport bounds for the identity
target cannot tend to zero.  No arithmetic estimate is supplied, and no
unconditional result is claimed.  The mathematical statements are
formalized in Lean~4 with Mathlib, using its statement of the Riemann
hypothesis; Appendix~A gives the scope.
\end{abstract}

\medskip
\noindent\textbf{Keywords.} Riemann hypothesis; zeros of the zeta function; conditional theorem; positive operators; self-dual trace confinement; Six Birds Theory; Lean formalization.

\section{Introduction}
\label{sec:intro}

The Riemann hypothesis \citep{Riemann1859} asserts that every
nontrivial zero \(\rho\) of the Riemann zeta function satisfies
\(\Real\rho=\frac12\).  Because the nontrivial zeros are the zeros of
the completed zeta function
\(\Lzeta(s)=\pi^{-s/2}\Gamma(s/2)\zeta(s)\) in the strip
\(0<\Real s<1\), the hypothesis can be restated as the vanishing of a
single nonnegative quantity,
\[
  \AZ=\sum_{\rho}m_\rho\Bigl(\Real\rho-\frac12\Bigr)^2
  \in[0,\infty],
\]
the sum running over the nontrivial zeros \(\rho\) with positive
weights \(m_\rho\) (for instance their multiplicities).  This
reformulation is elementary (\Cref{thm:rh:translation-T}); its
interest here is that it places the Riemann hypothesis inside the
setting of a companion paper \citep{TsiokosSDTC2026}.  There \(J\) is
an involution of a measure space \(X\), \(\psi\) is a strongly
measurable observable with values in a real Hilbert space that
intertwines \(J\) with a linear isometric involution, and its odd part
\(\psim\) is square-integrable and vanishes exactly on the fixed locus
of \(J\).  The Gram operator \(\AX=\int_X\psim\psim^*\) is then shown to
vanish---and the mass to lie on the fixed locus---as soon as it is
dominated by positive trace-class operators of vanishing trace; the
proof is a one-line trace squeeze.  The global square-integrability hypothesis would require \(\AZ<\infty\),
which is not known to hold unconditionally; the
applications below use either an abstract form of the theorem or
finite windows of zeros, where no such issue arises.  For the zeros of \(\zeta\) the relevant
involution is
\[
  \JL(s)=1-\bar s,
\]
reflection in the critical line, which maps zeros to zeros by the
functional equation and the reflection principle; its fixed locus is
the critical line, and the odd readout is \(\psimRH(s)=\Real s-\frac12\).

This paper examines that route carefully.  Its conclusions are as
follows.

\begin{enumerate}
  \item \emph{A conditional theorem.}  A named hypothesis
    \(\GamSDTCSelberg\) (\Cref{obl:rh:gamma-sdtc-selberg}) supplies
    positive domination records of vanishing trace for an auxiliary
    duality-confinement ledger, together with a bridge encoding, at each represented zero \(\rho\) of a
    shell (\Cref{sec:sat_sel_shell}), the displacement coordinate
    \(c(\rho)-\frac12\) in that ledger's odd readout.  Under
    this hypothesis every represented zero has coordinate
    \(c(\rho)=\frac12\) (\Cref{thm:rh:conditional}).  If,
    in addition, a map from coordinates to real numbers, sending
    \(\frac12\) to \(\frac12\), carries the represented zeros onto the
    real parts of all nontrivial zeros, then the Riemann hypothesis
    holds (\Cref{thm:rh:full-classical}).
  \item \emph{A calibration of the hypothesis.}  For every shell, the hypothesis holds exactly when every represented zero
    has coordinate \(\frac12\); a one-point ledger suffices
    (\Cref{prop:rh:gamma-calibration}).  On the bare shell, whose
    represented zeros are the nontrivial zeros with their real parts as
    coordinates, the hypothesis is therefore equivalent to the Riemann
    hypothesis.  The same is true of the operator form, with the actual zeros as
    objects, in which the duality-confinement theorem is applied to a
    sequence of finite windows of actual zeros,
    each closed under \(\JL\), that eventually contains every zero
    (\Cref{thm:rh:concrete-window-dc}).  The conditional theorem
    is therefore a reformulation of the Riemann hypothesis in the
    language of domination records, not evidence for it.
  \item \emph{Finite windows.}  On a finite window \(W\) of zeros
    closed under the functional equation, with weights invariant under
    it, the window sum \(\sum_{\rho\in W}m_\rho(\Real\rho-\frac12)^2\)
    is the squared norm of a vector that is invisible to the natural
    invariant probe (\Cref{sec:finite_windows}).  This identifies several ``payments''---upper bounds for the window
    sums, or for operators dominating them, that one might hope to
    establish independently.
    These facts are elementary.  Along windows that eventually contain
    every zero, each of several natural payment conditions is again
    equivalent to the Riemann hypothesis; and on a nonempty window, operator-level transport bounds for the
    identity target, of the form made precise in
    \Cref{prop:rh:window-obstructions}, cannot tend to zero.
\end{enumerate}

No unconditional statement about the zeros of \(\zeta\) is claimed.
Two features should be said at the outset.  First, no quantitative or prime-side input is used or supplied beyond
the classical zero-free half-plane.  The arguments use
only analytic facts about \(\zeta\)---the functional equation and
conjugation symmetry, countability of the zeros, the poles and non-vanishing of the Gamma factor, and the absence of zeros with \(\Real s\geq1\)---and
for any countable set of points in the strip with positive weights and
the same symmetries they prove the corresponding statements, with that
set's own critical-line property in place of the Riemann hypothesis.  Second, the functional-analytic content of the
companion theorem is not exercised: the conditional theorem uses an
abstract trace squeeze, and the concrete form uses operators on
\(\mathbb R\), that is, real numbers.  What the paper contributes is a
precise account of what a proof along these lines would have to supply,
and a demonstration that the obvious candidates are no easier than the
Riemann hypothesis itself.  The classical programs aimed at the Riemann
hypothesis---spectral interpretations \citep{BerryKeating1990},
trace formulas in noncommutative geometry \citep{Connes1999}, Hilbert
spaces of entire functions \citep{deBranges1968,deBranges1986}, and
Weil's positivity criterion \citep{Weil1952,Bombieri2000} and Li's
criterion \citep{Li1997}---are not used, and nothing here offers an
alternative to them.

The mathematical statements of
\Cref{sec:involution_and_ledger,sec:sat_sel_shell,sec:translation_theorem,sec:recognition_source,sec:landing_chain,sec:finite_windows}
are formalized in Lean~4 with the Mathlib library, and the Riemann
hypothesis in the conclusions is Mathlib's statement of it; the records
of \Cref{sec:rh-aor-instance} are formalized as lists of labels, and
the remarks are not formalized (\Cref{app:formalization}).  The
conditional theorem is stated for a \emph{shell}, a record whose
fields name the data an intended trace-formula realization would
provide (\Cref{sec:sat_sel_shell}); beyond the represented zeros and
their coordinates these data carry no analytic content.

\paragraph{Organization.}
\Cref{sec:framework} explains the framework vocabulary.
\Cref{sec:involution_and_ledger} defines the involution, the readout
and the zero ledgers; \Cref{sec:sat_sel_shell} describes the shell;
\Cref{sec:translation_theorem} proves that the ledger vanishes exactly
when the zeros lie on the critical line.
\Cref{sec:recognition_source,sec:landing_chain} state the hypothesis
\(\GamSDTCSelberg\), prove the conditional theorems and their
calibration.  \Cref{sec:rh-aor-instance} records a classification of
the argument's hypotheses in the audit vocabulary of a further
companion paper.  \Cref{sec:finite_windows} treats finite windows.
\Cref{sec:scope_and_nonclaims,sec:discussion,sec:conclusion} discuss
scope and context.
 \section{Framework}
\label{sec:framework}

This section explains the framework terms used in the paper.  No
proof depends on them, and a reader interested only in the
mathematics may skip to \Cref{sec:involution_and_ledger}.

The author's Six Birds framework \citep{TsiokosFoundationsI2026}
describes a mathematical theory as the collection of fixed points of
an idempotent completion rule; later papers add a discipline for which
candidate objects are admissible \citep{TsiokosFoundationsII2026} and
a finite vocabulary for what an observer can see and certify
\citep{TsiokosFoundationsIII2026}.  Within it, a \emph{structural law}
is a named hypothesis asserting that a specified class of structures
carries certain data, and the data it asserts are called its
\emph{recognition source}.  The point of the naming is bookkeeping:
every result that uses the hypothesis can be seen to use it.

The structural law relevant here is Self-Dual Trace Confinement
(SDTC) of the companion paper \citep{TsiokosSDTC2026}, in the setting
recalled in \Cref{sec:intro} (a strongly measurable observable in a
real Hilbert space intertwining an involution with a linear isometric
involution): in a specified class of closures with a genuine involutive self-duality, observed
through a separating odd readout, the structure supplies positive trace-class operators \(B_n\) with
\(\AX\preceq B_n\) and \(\trace B_n\to0\), where \(\AX\) is the Gram
operator of the square-integrable odd readout.
That paper proves the corresponding confinement theorem and its
converse, and shows that for a single ledger with a separating readout
the hypothesis is equivalent to confinement.  The present paper
specializes the hypothesis to the zeros of the zeta function, where it
is called \(\GamSDTCSelberg\): the subscript records that its intended
source is a trace formula of Selberg type, which is not constructed
here.  Consistently with the companion paper, \(\GamSDTCSelberg\) turns out to be equivalent to the Riemann hypothesis
on the bare shell of all nontrivial zeros with their real parts
(\Cref{sec:recognition_source,sec:landing_chain}).

Two further companion papers supply vocabulary used in
\Cref{sec:rh-aor-instance,sec:finite_windows}: an audit calculus for
classifying the hypotheses of an argument \citep{TsiokosAOR2026}, and
a residual operator measuring what a probe fails to see
\citep{TsiokosXiCompanion2026}.  Only elementary special cases of the
latter are used, and they are defined where needed.
 \section{Involution and ledger}
\label{sec:involution_and_ledger}

This section introduces the functional-equation involution, its odd
readout, and the two zero ledgers.  Everything here is elementary; the
point is to place the critical line in the pattern of an involution,
its fixed locus, and a readout that changes sign.

\begin{definition}[Functional-equation involution \(\JL\)]
\label{def:rh:fe-involution}
The functional-equation involution is the map
\(\JL:\mathbb C\to\mathbb C\),
\[
  \JL(s)=1-\bar s .
\]
If \(s=x+iy\), then \(\JL(s)=(1-x)+iy\): the map is reflection in the
vertical line \(\Real s=\frac12\).
\end{definition}

Directly from the formula, \(\JL(\JL(s))=1-\overline{1-\bar s}=s\), and
\(\JL(s)=s\) if and only if \(\Real s=\frac12\).  Thus \(\JL\) is an
involution whose fixed locus is the critical line.  The completed zeta
function satisfies \(\Lzeta(1-s)=\Lzeta(s)\) and
\(\Lzeta(\bar s)=\overline{\Lzeta(s)}\)
\citep{TitchmarshHeathBrown1986,Edwards1974}, so \(\JL\) maps zeros of
\(\Lzeta\) to zeros of \(\Lzeta\), preserving the strip
\(0<\Real s<1\).

\begin{definition}[Anti-invariant readout \(\psimRH\)]
\label{def:rh:psi-minus-rh}
The anti-invariant readout of \(\JL\) is the signed horizontal
displacement from the critical line,
\[
  \psimRH(s)=\Real s-\frac12 .
\]
\end{definition}

Since \(\Real\JL(s)=1-\Real s\), the readout is odd:
\(\psimRH(\JL(s))=-\psimRH(s)\).  It vanishes exactly on the fixed locus
of \(\JL\), so it is separating in the sense of the companion paper.

\begin{definition}[Nontrivial-zero ledger]
\label{def:rh:nontrivial-zero-ledger}
The \emph{nontrivial zeros} are the zeros \(\rho\) of \(\Lzeta\) with
\(0<\Real\rho<1\); they are exactly the zeros of \(\zeta\) in that
strip.  A \emph{nontrivial-zero ledger} is the set \(\Znt\) of
nontrivial zeros together with positive integer weights
\(m_\rho\) \((\rho\in\Znt)\).  The natural choice of \(m_\rho\) is the
multiplicity of \(\rho\) as a zero of \(\Lzeta\); every result below holds for arbitrary positive integer weights, and the formal development uses
unit weights where a choice is needed.
\end{definition}

\begin{definition}[Anti-invariant zero ledger]
\label{def:rh:anti-invariant-zero-ledger}
The \emph{anti-invariant zero ledger} is the extended nonnegative sum
\[
  \AZ=\sum_{\rho\in\Znt}m_\rho\,\bigl|\psimRH(\rho)\bigr|^2
     =\sum_{\rho\in\Znt}m_\rho\Bigl(\Real\rho-\frac12\Bigr)^2
  \in[0,\infty],
\]
defined as the supremum of its finite partial sums.
\end{definition}

No convergence assumption is needed: the value \(+\infty\) is allowed,
and it is only the vanishing of \(\AZ\) that matters.  Because every
summand is nonnegative and every weight is positive, \(\AZ=0\) exactly
when every nontrivial zero lies on the critical line;
\Cref{sec:translation_theorem} states this as Theorem~T.  For a finite
set \(W\subseteq\Znt\), the \emph{window sum}
\(A_W=\sum_{\rho\in W}m_\rho(\Real\rho-\frac12)^2\) is a real number
bounded by \(\AZ\); windows are the subject of
\Cref{sec:finite_windows}.
 \section{The saturated Selberg trace shell}
\label{sec:sat_sel_shell}

The conditional theorem of \Cref{sec:landing_chain} is stated for a
record that collects, under one name, the data an intended
trace-formula realization of the zeros would provide.  We call it the
saturated completed Selberg trace shell and write \(\Sel\).  The name records the intended realization, through a trace formula of
Selberg type \citep{Selberg1956,Hejhal1976,Iwaniec2002}.  Spectral realizations of the zeros of \(\zeta\) exist---Connes' trace
formula \citep{Connes1999} and Meyer's spectral interpretation
\citep{Meyer2005}, which realizes Weil's explicit formula
\citep{Weil1952} as a trace formula---but none is known to come with a self-adjoint operator \(H\) such that
every nontrivial zero is \(\frac12+i\lambda\) with \(\lambda\) in the spectrum
of \(H\), which would prove the Riemann hypothesis.  (For the Selberg zeta function of a compact hyperbolic
surface, self-adjointness of the Laplacian places the zeros coming from
eigenvalues \(\lambda\geq\frac14\) on the critical line, while eigenvalues
in \((0,\frac14)\) give finitely many real exceptional zeros.)  Nothing of this kind is constructed or used here.

\begin{definition}[Saturated completed Selberg trace shell]
\label{def:rh:sat-sel-shell}
A \emph{saturated completed Selberg trace shell} \(\Sel\) consists of
the following data:
\begin{enumerate}[label=\textup{(\arabic*)}]
  \item sets \(H_L\), \(\Itr\), \(\mathrm{EQ}_{\mathrm{tr}}\) and
        \(\mathrm{EM}_{\mathrm{zero}}\), naming a history carrier, a
        trace instrument, trace observables and verifier events;
  \item sets \(Q_L\) and \(M_L\) with a map \(\pi_L:M_L\to Q_L\),
        naming a current and a predictive quotient and their
        comparison;
  \item a \emph{coordinate system}: a set \(R\) with elements \(0\) and
        \(\frac12\) and operations \(x\mapsto-x\), \((x,y)\mapsto x-y\) and
        \(x\mapsto1-x\), such that \(x\mapsto1-x\) is an involution whose
        only fixed point is \(\frac12\), \(x-\frac12=0\) if and only if
        \(x=\frac12\), and \((1-x)-\frac12=-(x-\frac12)\); the
        \emph{points} of the shell are pairs \(s=(x(s),y(s))\in R^2\), and the \emph{coordinate} of a point \(s\) is \(x(s)\);
  \item the functional-equation involution
        \(\JL(s)=(1-x(s),y(s))\) on points,
        whose fixed points are those with \(x(s)=\frac12\),
        a set \(V\) with a subset of \emph{zero values}, and a function
        \(\Lambda_L\) from points to \(V\) (the intended completed zeta
        function);
  \item a zero ledger: an element \(1\in R\), a relation \(<\) on \(R\)
        (no laws are imposed on it), a set of \emph{represented zeros}
        with an injective map \(\rho\mapsto\) point, and positive integer
        weights \(m_\rho\); we write \(c(\rho)\) for the coordinate of the
        point of \(\rho\).  It is required that every represented zero satisfies
        \(0<c(\rho)<1\) and \(\Lambda_L(\rho)\) is a zero
        value, and every point \(s\) with \(0<x(s)<1\) and
        \(\Lambda_L(s)\) a zero value is represented;
  \item a function \(q:R\to[0,\infty]\) vanishing exactly at \(0\) (the
        intended square), and the anti-invariant ledger
        \(\AZ=\sum_\rho m_\rho\,q(c(\rho)-\frac12)\) of the
        represented zeros, an extended nonnegative sum;
  \item a visibility relation \(\mathrm{Vis}_L\) between histories and
        instruments, and an audit marker \(\mathrm{Audit}_L\).
\end{enumerate}
\end{definition}

Of these data, the proofs in this paper use only the represented zeros,
their weights and coordinates, the law ``\(x-\frac12=0\) if and only if
\(x=\frac12\)'', and \(q\).  The remaining laws and data describe the
intended realization and serve to specify the bare shell below; the
operation \(x\mapsto1-x\) is a formal symbol subject only to the laws
listed.  Items (1), (2) and (7) are names: the shell imposes no laws on them, and in particular
it does not assert that the trace instrument satisfies a trace formula,
that the quotients have any universal property, or that the shell is
admissible in the sense of the framework's formation rules
\citep{TsiokosFoundationsII2026}; the audit marker records a status,
not a proof of admissibility.  A shell is useful precisely because it
separates the zeros it represents from the actual zero set: the conditional theorem speaks about represented zeros and their
coordinates, and \Cref{thm:rh:full-classical} adds the hypothesis that a
map from coordinates to the reals, sending \(\frac12\) to \(\frac12\),
carries them onto the real parts of all nontrivial zeros.  One shell is always available: the \emph{bare shell}, with
\(R=\mathbb R\) and its usual operations, element \(1\) and order, points
\(\mathbb C\), \(V=\mathbb C\) with zero value \(0\), \(\Lambda_L=\Lzeta\),
represented zeros exactly the nontrivial zeros of
\Cref{def:rh:nontrivial-zero-ledger} with unit weights, \(q(x)=x^2\),
and trivial documentary fields (1), (2) and (7).  For it
\(c(\rho)=\Real\rho\) and \(\AZ\) is the ledger of
\Cref{def:rh:anti-invariant-zero-ledger}.
 \section{Translation theorem T}
\label{sec:translation_theorem}

The next theorem translates the vanishing of the anti-invariant ledger
into the critical-line statement.  It is a positivity argument.  We
state it for the nontrivial-zero ledger of
\Cref{def:rh:nontrivial-zero-ledger}; the identical argument applies to
the ledger of represented zeros of any shell, which is how it is used
in \Cref{sec:landing_chain}.

\begin{theorem}[Translation theorem T]
\label{thm:rh:translation-T}
For any positive weights \(m_\rho\),
\[
  \AZ=0
  \quad\Longleftrightarrow\quad
  \Real\rho=\frac12\ \text{ for every }\rho\in\Znt .
\]
Consequently \(\AZ=0\) if and only if the Riemann hypothesis holds.
\end{theorem}

\begin{proof}
The two directions are
\Cref{thm:rh:translation-T-forward,thm:rh:translation-T-reverse}.  The
nontrivial zeros are exactly the zeros of \(\zeta\) in the strip
\(0<\Real s<1\), because \(\pi^{-s/2}\Gamma(s/2)\) has neither zeros
nor poles there; and every zero of \(\zeta\) other than the trivial zeros \(-2,-4,\ldots\)
lies in that strip.  For the latter, the classical facts used are
\citep{TitchmarshHeathBrown1986}: \(\zeta\) has no zeros with
\(\Real s\geq1\); \(\zeta(0)=-\frac12\); \(\Lzeta(s)=\Lzeta(1-s)\); and
\(\pi^{-s/2}\Gamma(s/2)\) never vanishes and is finite except at
\(s=0,-2,-4,\ldots\).  If \(\zeta(s)=0\) with \(\Real s\leq0\), then
\(s\neq0\), and if \(s\) is not a trivial zero the factor at \(s\) is
finite and nonzero, so \(\Lzeta(s)=0\), hence \(\Lzeta(1-s)=0\); since
\(\Real(1-s)\geq1\) and \(1-s\neq1\), the factor at \(1-s\) is finite and
nonzero, so \(\zeta(1-s)=0\), a contradiction.  Hence the critical-line
statement on \(\Znt\) is the Riemann hypothesis.
\end{proof}

\begin{theorem}[Translation theorem T, forward direction]
\label{thm:rh:translation-T-forward}
If \(\AZ=0\), then \(\Real\rho=\frac12\) for every \(\rho\in\Znt\).
\end{theorem}

\begin{proof}
Each summand \(m_\rho(\Real\rho-\frac12)^2\) is nonnegative and is at
most the extended sum \(\AZ\).  If \(\AZ=0\), each summand is therefore
zero, and since \(m_\rho>0\) we get \((\Real\rho-\frac12)^2=0\).
\end{proof}

\begin{theorem}[Translation theorem T, reverse direction]
\label{thm:rh:translation-T-reverse}
If \(\Real\rho=\frac12\) for every \(\rho\in\Znt\), then \(\AZ=0\).
\end{theorem}

\begin{proof}
Every summand vanishes, so every partial sum vanishes, and so does
their supremum.
\end{proof}

The equivalence between \(\AZ=0\) and the critical-line statement uses
nothing about the zeros beyond positivity of the weights; only its
identification with the Riemann hypothesis uses the analytic facts
recalled in the proof.  Its role is to let the rest of the paper work with the single
quantity \(\AZ\), or with its finite window sums, instead of with the
zero set.
 \section{The recognition source}
\label{sec:recognition_source}

By Theorem~T, the Riemann hypothesis is the statement \(\AZ=0\).  The
duality-confinement route asks for this vanishing to come from
domination records of vanishing trace.  The following hypothesis
states what the conditional theorem of \Cref{sec:landing_chain}
consumes; the formal record carries in addition some structure that no
proof uses (weights, a response involution, a quantitative separation
modulus), listed in \Cref{app:formalization}.  It is phrased with an auxiliary
ledger in the abstract form of the duality-confinement theorem
\citep[Remark~7.2]{TsiokosSDTC2026}, because the formal statement
allows the domination records to live in any ordered set with a trace,
not only among operators.

\begin{obligation}[The recognition source \(\GamSDTCSelberg\)]
\label{obl:rh:gamma-sdtc-selberg}
Let \(\Sel\) be a shell.  We say that the recognition source
\(\GamSDTCSelberg\) \emph{holds for} \(\Sel\) if there exist the
following data.
\begin{enumerate}[label=\textup{(\roman*)}]
  \item \emph{An auxiliary ledger.}  A finite set \(X\) with an involution \(J\) (the map
    \(\rho\mapsto x_\rho\) of item (iv) need not be injective), a set of \emph{visible} objects, a response set
    \(Y\) with a zero element, and a readout
    \(\psim:X\to Y\) that is separating on visible objects:
    for visible \(x\), \(\psim(x)=0\) if and only if \(Jx=x\).
  \item \emph{An ordered set with a trace.}  A set \(C\) with a
    relation \(\preceq\), a subset of positive elements, a zero element,
    a distinguished positive element \(\AX\), and a map
    \(\tau:C\to\mathbb R\) that is monotone for \(\preceq\) and
    nonnegative on positive elements, with \(\tau(\AX)=0\) only if
    \(\AX\) equals the zero element.
  \item[(ii\('\))] \emph{A link between the two.}  If \(\AX\) is the zero
    element, then \(\psim(x)=0\) at every visible \(x\).
  \item \emph{Domination records of vanishing trace.}  Positive
    elements \(B_n\in C\) with \(\AX\preceq B_n\) for all \(n\) and
    \(\tau(B_n)\to0\).
  \item \emph{A bridge to the zeros.}  For each represented zero
    \(\rho\) of \(\Sel\), a visible object \(x_\rho\in X\), and a map
    \(e:R\to Y\) on the shell's coordinates with \(e(r)=0\) only if
    \(r=0\), such that
    \[
      \psim(x_\rho)=e\Bigl(c(\rho)-\frac12\Bigr).
    \]
\end{enumerate}
\end{obligation}

Items (i)--(iii) are an instance of the duality-confinement setting in
its abstract form: the domination records and the trace squeeze force
\(\AX\) to vanish, and item (ii\('\)) then gives \(\psim=0\) at every visible
object.  None of this involves the zeta function, and neither \(\AZ\) nor \(q\)
appears in the hypothesis; the zeros enter only through item (iv), which asserts that the auxiliary readout at the
object attached to \(\rho\) encodes the displacement coordinate
\(c(\rho)-\frac12\).  In this abstract form the link between
\(\AX\) and \(\psim\) is part of the data (item (ii\('\))) rather than a
consequence of a construction, and the next proposition shows that the
hypothesis is exactly as strong as the critical-line statement for the
represented zeros.  The operator form, in which \(\AX\) is the Gram
operator of the readout, is \Cref{thm:rh:concrete-window-dc}.  Nothing in the
hypothesis requires the auxiliary ledger to be constructed from
zeta-function data, and its intended provenance---an instance of the
companion paper's structural law obtained from a trace formula---is
not established here.

The next proposition measures the strength of the hypothesis.

\begin{proposition}[Calibration of the recognition source]
\label{prop:rh:gamma-calibration}
For every shell \(\Sel\), \(\GamSDTCSelberg\) holds for \(\Sel\) if and
only if \(c(\rho)=\frac12\) for every represented zero
\(\rho\).  In particular, \(\GamSDTCSelberg\) holds for the bare shell if and only
if the Riemann hypothesis holds.
\end{proposition}

\begin{proof}
If \(\GamSDTCSelberg\) holds, \Cref{thm:rh:conditional} below gives
\(c(\rho)=\frac12\) for every represented zero.
Conversely, suppose \(c(\rho)=\frac12\) for all represented
zeros.  Take \(X\) a single visible point with \(J\) the identity,
\(Y=\mathbb R\), \(\psim=0\), \(C=\mathbb R\) with the usual order and positive elements
\([0,\infty)\),
\(\AX=0\), \(B_n=0\), \(\tau\) the identity, \(x_\rho\) the point, and \(e:R\to\mathbb R\) with \(e(0)=0\) and \(e(r)=1\) for \(r\neq0\)
(a classical case distinction on the abstract set \(R\)).  Items (i)--(iii) hold trivially, \(e(r)=0\) only for \(r=0\), and item (iv)
holds because \(c(\rho)-\frac12=0\) by the coordinate
laws, so \(e(c(\rho)-\frac12)=0=\psim(x_\rho)\).  On the
bare shell, \(c(\rho)=\Real\rho\) and the represented zeros
are the nontrivial zeros, so the condition is the Riemann hypothesis
(\Cref{thm:rh:translation-T}).
\end{proof}

Thus the recognition source, as a type of data, is exactly as strong as
the critical-line statement it is used to prove, and on the bare shell
establishing it is no easier than proving the Riemann
hypothesis.  Its intended content is not its logical
strength but its provenance: the conditional theorem becomes a proof
only if the data are produced by an independent argument, for instance
from a trace formula, in which item (iv) is a consequence rather than
an assumption.  \Cref{thm:rh:concrete-window-dc} shows that the same
calibration holds for a non-degenerate version built from the actual
zeros.
 \section{The conditional landing chain}
\label{sec:landing_chain}

This section proves the conditional theorems and the concrete
calibration.  The reader may keep the bare shell in mind throughout: for
it \(c(\rho)=\Real\rho\), the represented zeros are all nontrivial
zeros, and \Cref{thm:rh:conditional,thm:rh:full-classical} say the same
thing.  The duality-confinement theorem is used in the abstract
form of the companion paper \citep[Remark~7.2]{TsiokosSDTC2026}: if
\(a\) is a positive element of an ordered set, \(\tau\) is a
nonnegative monotone real trace with \(\tau(a)=0\) only if \(a=0\), and
\(a\preceq b_n\) with \(\tau(b_n)\to0\), then \(a=0\), because
\(0\leq\tau(a)\leq\tau(b_n)\to0\).

\begin{theorem}[Duality-confinement theorem applied to \(\Sel\)]
\label{thm:rh:dc-master-applied}
Let \(\Sel\) be a shell for which \(\GamSDTCSelberg\) holds
(\Cref{obl:rh:gamma-sdtc-selberg}).  Then the
anti-invariant ledger of the represented zeros vanishes: \(\AZ=0\).
\end{theorem}

\begin{proof}
By items (ii) and (iii) and the abstract squeeze, \(\AX\) is the zero
element; by item (ii\('\)), \(\psim(x)=0\) at every visible object \(x\).
For a represented zero \(\rho\), the object \(x_\rho\) is visible, so
by item (iv) \(e(c(\rho)-\frac12)=\psim(x_\rho)=0\), and
hence \(c(\rho)=\frac12\).  (Separation, item (i), gives the equivalent
statement that \(x_\rho\) is fixed by \(J\); it is not needed for the
conclusion.)  Every summand of \(\AZ\) therefore vanishes, and
\(\AZ=0\) by \Cref{thm:rh:translation-T-reverse} applied to the represented zeros.
\end{proof}

The companion theorem enters this proof only through the abstract
squeeze; the link from \(\AX\) to the readout is item (ii\('\)) and the
link to the zeros is item (iv), both of which are assumed.

\begin{theorem}[Conditional landing chain]
\label{thm:rh:conditional}
Let \(\Sel\) be a shell for which \(\GamSDTCSelberg\) holds.  Then every represented zero \(\rho\) of \(\Sel\)
satisfies \(c(\rho)=\frac12\).  On the bare shell this says
\(\Real\rho=\frac12\) for every nontrivial zero \(\rho\).
\end{theorem}

\begin{proof}
\Cref{thm:rh:dc-master-applied} gives \(\AZ=0\), and the forward
direction of Theorem~T (\Cref{thm:rh:translation-T-forward}), applied
to the represented zeros, gives the conclusion.
\end{proof}

\begin{theorem}[Full classical conditional Riemann hypothesis]
\label{thm:rh:full-classical}
Let \(\Sel\) be a shell with coordinate system \(R\), and suppose that
\(\GamSDTCSelberg\) holds for \(\Sel\) and that there is a
\emph{classical-zero identification}: a map \(t:R\to\mathbb R\) with \(t(\frac12)=\frac12\)
such that for every nontrivial zero \(s\) of \(\zeta\) there is a
represented zero \(\rho\) with \(t(c(\rho))=\Real s\).
Then the Riemann hypothesis holds.  For the bare shell, \(t\) is the
identity.
\end{theorem}

\begin{proof}
By \Cref{thm:rh:conditional}, \(c(\rho)=\frac12\) for every
represented zero, so \(\Real s=t(\frac12)=\frac12\) for every nontrivial
zero \(s\).  The Riemann hypothesis follows, because every zero of \(\zeta\) other
than the trivial zeros lies in the strip \(0<\Real s<1\) (proof of
\Cref{thm:rh:translation-T}).
\end{proof}

Taken together with \Cref{prop:rh:gamma-calibration}, these theorems
say that the Riemann hypothesis is equivalent to the statement that \(\GamSDTCSelberg\) holds for the bare shell.  The recognition source
used there is degenerate: a single object carries the readout.  The
next theorem applies the operator form of the duality-confinement
theorem to the actual zeros, with the zeros themselves as objects.
  Let
\((W_n)_{n\geq1}\) be finite sets of nontrivial zeros, each closed under
\(\JL\), such that every nontrivial zero lies in \(W_n\) for all large
\(n\); such a sequence exists because the nontrivial zeros are countable and
each orbit of \(\JL\) has one or two elements (take \(W_n\) to be the
union of the orbits of the first \(n\) zeros in an enumeration).  On
\(W_n\) take the finite weighted ledger of the companion
paper with unit weights, response space \(\mathbb R\), response
involution \(v\mapsto-v\), involution \(\JL\) and readout
\(\psimRH(\rho)=\Real\rho-\frac12\).  Its operator is multiplication by the unit-weight window sum
\(A_{W_n}=\sum_{\rho\in W_n}(\Real\rho-\frac12)^2\), and its readout is
separating: it vanishes exactly at the fixed points of \(\JL\).  The
response space is one-dimensional, so the operators are scalars and the
theorem is the simplest case of the duality-confinement theorem; what
it adds to \Cref{prop:rh:gamma-calibration} is that the objects are the
actual zeros and the bridge is an identity rather than an
assumption.

\begin{theorem}[Concrete window form]
\label{thm:rh:concrete-window-dc}
With the windows \(W_n\) above, the following are equivalent:
\begin{enumerate}[label=\textup{(\alph*)}]
  \item for every \(n\) there are real numbers \(b_{n,k}\geq A_{W_n}\)
    with \(b_{n,k}\to0\) as \(k\to\infty\);
  \item there are real numbers \(b_n\geq A_{W_n}\) with \(b_n\to0\);
  \item the Riemann hypothesis.
\end{enumerate}
In (a) and (b) the numbers \(b\) are the traces of positive operators
on \(\mathbb R\) dominating the window operators, so (a) is the
hypothesis of the duality-confinement theorem for each window, and the
implication (a)\(\Rightarrow\)(c) is proved by that theorem.
\end{theorem}

\begin{proof}
(a)\(\Rightarrow\)(c): for each \(n\), \(0\leq A_{W_n}\leq b_{n,k}\to0\), so
\(A_{W_n}=0\); this is the duality-confinement theorem for this finite
weighted ledger \citep[Theorem~7.1]{TsiokosSDTC2026}, and since the readout
is separating every zero in \(W_n\) is fixed by \(\JL\), that is, lies on
the critical line; every nontrivial zero lies in some
window, and \Cref{thm:rh:translation-T} gives the Riemann hypothesis.
(b)\(\Rightarrow\)(c): if some nontrivial zero \(\rho_0\) were off the
line, it would lie in \(W_n\) for all large \(n\), so
\(b_n\geq A_{W_n}\geq(\Real\rho_0-\frac12)^2>0\) for all large \(n\),
contradicting \(b_n\to0\).  (c)\(\Rightarrow\)(a) and
(c)\(\Rightarrow\)(b): under the Riemann hypothesis every window sum is
zero, and the zero numbers dominate it.
\end{proof}

Unwound, condition (a) says that every window sum is zero and
condition (b) that the window sums tend to zero; the theorem makes
explicit that, for the actual zeros, domination records of vanishing
trace are nothing more than such upper bounds, and that they exist
exactly when the window sums vanish.  A
proof of the Riemann hypothesis along this route therefore needs an
independent source of such bounds; \Cref{sec:finite_windows} examines
some candidate sources.
 \section{A classification of the hypotheses}
\label{sec:rh-aor-instance}

A companion paper on audited operational realisability (AOR)
\citep{TsiokosAOR2026} provides a vocabulary for classifying the
ingredients of an argument: each ingredient receives a \emph{residual
type} describing what kind of gap it closes (a missing source, a
transport between settings, a presentation choice, a meta-level
condition, and so on) and a \emph{status} describing how the gap is
closed (by construction, by a supplied bridge, by another paper, or
not at all).  This section records such a classification for the
argument of \Cref{sec:landing_chain}.  It is bookkeeping: it proves nothing about the zeros, can be skipped
by a reader interested only in the mathematics, and does not establish
membership of the
argument in the stronger, proof-carrying notion of realisability of
that paper.

\begin{definition}[The AOR instance object attached to \(\Sel\)]
\label{def:rh:aor-sel-instance}
Let \(\mathfrak S_{\mathrm{RH}}\) be the scope consisting of the shell
\(\Sel\), the objects of \Cref{sec:involution_and_ledger}, the
recognition source \(\GamSDTCSelberg\), the duality-confinement theorem
and Theorem~T.  The AOR instance object
\(C_{\mathrm{RH}}(\Sel,\GamSDTCSelberg)\) is the list of the following
entries, each with a residual type and a status:
the carrier (the zero side of \(\Sel\) with its quotients); the
observations (\(\Znt\), \(\psimRH\), \(\AZ\), \(\mathrm{Vis}_L\)); the
routes (\(\JL\), \(\pi_L\), and the bridge from the auxiliary ledger to
the shell); the sources (\(\mathrm{Audit}_L\), \(\GamSDTCSelberg\) and
its records \(B_n\)); the constraints (positivity, order, trace
monotonicity, domination, the trace limit); the records \labelcref{thm:rh:aor-mechanical-records,thm:rh:aor-recognition-discharge,thm:rh:aor-gamma-bridge-discharge,thm:rh:aor-auditL-discharge,thm:rh:aor-dc-master-import-discharge,thm:rh:aor-real-coordinate-discharge,thm:rh:aor-translation-interface-discharge}
below; and the non-claims of \Cref{sec:scope_and_nonclaims}.
\end{definition}

Each record below is a short list of \emph{atoms}.  An atom consists of
a \emph{primary} residual type (the kind of gap the entry closes), a
list of \emph{forced secondary} types (further kinds of gap that the
entry is declared to bring with it, each of which must be recorded in
its own right), and a status.  The statuses used are: \emph{by construction} (a definition),
\emph{zero} (proved, so that no gap remains), \emph{bridged} (supplied
as a hypothesis), \emph{asymptotically budgeted} (a limit supplied as a
hypothesis), \emph{proved in another paper}, \emph{outside scope}, and
\emph{non-claim}.  \Cref{thm:rh:aor-instance} below is a consistency
check on these lists: it says that every kind of gap an entry declares
is itself recorded by some entry, and nothing more.

\begin{classification}[Mechanical records]
\label{thm:rh:aor-mechanical-records}
The mechanical record has six atoms, one for each of
\Cref{def:rh:fe-involution,def:rh:psi-minus-rh,def:rh:nontrivial-zero-ledger,def:rh:anti-invariant-zero-ledger,def:rh:sat-sel-shell}
and \Cref{thm:rh:translation-T}.  The five definition atoms have
primary types role (involution), role (readout), presentation (zero
ledger), presentation (anti-invariant ledger) and meta (shell), each
with status by construction and no forced secondaries; the atom for
Theorem~T has primary type target, forced secondary transport, and
status zero.  None uses recognition-source content.
\end{classification}

\begin{classification}[Recognition-source record]
\label{thm:rh:aor-recognition-discharge}
\(\GamSDTCSelberg\) is recorded by an atom with primary type source,
forced secondaries target, role and limit, and status bridged: it is
supplied, not derived.  Three further atoms record the secondaries:
target and role, bridged, and limit, asymptotically budgeted.
\end{classification}

\begin{classification}[Bridge-field record]
\label{thm:rh:aor-gamma-bridge-discharge}
The bridge between the auxiliary ledger and the shell (item (iv) of
\Cref{obl:rh:gamma-sdtc-selberg} with the readout-coherence fields) is
recorded by transport atoms---one with forced secondaries role and
target and status zero, two bridged---and bridged atoms for role and
target.
\end{classification}

\begin{classification}[Admissibility-audit record]
\label{thm:rh:aor-auditL-discharge}
The audit marker \(\mathrm{Audit}_L\) is recorded by one atom of
primary type meta and status bridged.  Admissibility of an intended
trace-formula realization under the framework's formation rules is not
established.
\end{classification}

\begin{classification}[Cross-paper record]
\label{thm:rh:aor-dc-master-import-discharge}
The use of the duality-confinement theorem of \citet{TsiokosSDTC2026}
is recorded by an atom of primary type transport with forced
secondaries source, role and target and status proved in another
paper, and bridged atoms for source, role and target.
\end{classification}

\begin{classification}[Real-coordinate record]
\label{thm:rh:aor-real-coordinate-discharge}
The abstract real coordinate of a general shell is recorded by two
atoms of primary type presentation, with statuses outside scope and
non-claim.  Its identification with \(\Real\rho\) is the separate
hypothesis of \Cref{thm:rh:full-classical}; for the bare shell it holds
by construction.
\end{classification}

\begin{classification}[Translation record]
\label{thm:rh:aor-translation-interface-discharge}
Theorem~T is recorded by an atom of primary type target with forced
secondary transport, and an atom of primary type transport, both with
status zero.
\end{classification}

\begin{theorem}[Local register completeness]
\label{thm:rh:aor-instance}
In the concatenation \(C_{\mathrm{RH}}(\Sel,\GamSDTCSelberg)\) of
\Cref{thm:rh:aor-mechanical-records,thm:rh:aor-recognition-discharge,thm:rh:aor-gamma-bridge-discharge,thm:rh:aor-auditL-discharge,thm:rh:aor-dc-master-import-discharge,thm:rh:aor-real-coordinate-discharge,thm:rh:aor-translation-interface-discharge},
the list of non-claims is nonempty, and every forced secondary type of
every atom occurs as the primary type of some atom.  In the terminology of \citet{TsiokosAOR2026}, the register
is locally complete on \(\mathfrak S_{\mathrm{RH}}\).
\end{theorem}

\begin{proof}
The non-claim list is nonempty by construction.  The
forced secondary types that occur are transport, target, role,
limit and source.  Each occurs as a primary type: transport in the bridge-field record,
target, role and limit in the recognition-source record, and source in
the recognition-source and cross-paper records.
\end{proof}

\begin{remark}[Partial status of the classification]
\label{rem:rh:aor-partial-status}
\Cref{thm:rh:aor-instance} is a statement about a finite list of
labels; it uses no domination record, bridge or trace limit, and it
does not imply the critical-line statement.  The proof-carrying notion
of realisability of \citet{TsiokosAOR2026} requires stored proofs for
each obligation and quantitative payment hypotheses, which are not
supplied here; membership of the argument in that notion remains open.
\end{remark}
 \section{Finite windows}
\label{sec:finite_windows}

By \Cref{thm:rh:concrete-window-dc}, a proof along the duality-confinement
route needs upper bounds tending to zero for the window sums \(A_W\).
This section records some elementary facts about natural candidates
for such bounds.  The main point is \Cref{prop:rh:window-identities}(b):
the symmetric probe \(L_W(x)=\sum x_\rho\) sees nothing of the odd
displacement vector, by symmetry alone, so its residual charge is the
whole window sum; the equivalences of
\Cref{prop:rh:window-equivalences} are then tautologies forced by this
choice of probe, and carry no information about \(\zeta\).  Only the oddness of the displacement under the
functional-equation map is used here; no fixed locus is involved.

\paragraph{Setting.}
Throughout, \(\sigma(\rho)=1-\rho\) is the functional-equation map on
the nontrivial zeros (not to be confused with the reflection \(\JL\) in
the critical line), a finite set \(W\) of nontrivial
zeros is \emph{reflected} if \(\sigma(W)=W\), and the positive integer
weights satisfy \(m_{\sigma(\rho)}=m_\rho\).  A sequence \((W_n)\) of
finite sets is \emph{exhaustive} if every nontrivial zero lies in
\(W_n\) for all large \(n\); exhaustive sequences of reflected windows
exist because the nontrivial zeros are countable and each orbit of
\(\sigma\) has at most two elements.  These windows differ from those of
\Cref{thm:rh:concrete-window-dc}, which are closed under
\(\JL(\rho)=1-\bar\rho\) and use unit weights; windows closed under both
reflections also exist, since the two reflections commute and each
orbit of the group they generate has at most four elements.  On the Euclidean
space \(\mathbb R^W\) put
\[
  v_W(\rho)=\sqrt{m_\rho}\Bigl(\Real\rho-\frac12\Bigr),
  \qquad
  L_W(x)=\sum_{\rho\in W}x_\rho ,
\]
the \emph{displacement vector} and the \emph{invariant probe}.  For a linear map \(L\) from \(\mathbb R^W\) to a finite-dimensional
real inner product space, let \(L^+\) be its Moore--Penrose pseudoinverse and
\[
  \Xi(L)=I-L^+L ,
\]
the orthogonal projection onto \(\ker L\); \(L^+L\) is the orthogonal
projection onto \((\ker L)^\perp\).  For a vector \(v\), the
\emph{residual charge} \(\langle v,\Xi(L)v\rangle=\|\Xi(L)v\|^2\) measures
the part of \(v\) that the probe \(L\) does not see.  We call \(LL^*\) the \emph{native currency} of \(L\).  (These are the
identity-audit cases of the residual, native currency and decoder
\(L^+\) of the companion paper \citep{TsiokosXiCompanion2026}.)  A \emph{payment} for a window sum is an upper bound for it,
or for a quadratic form dominating it, that one hopes to derive
independently.  Operator inequalities \(S\preceq T\) are in the sense of
quadratic forms.

\begin{proposition}[Window identities]
\label{prop:rh:window-identities}
Let \(W\) be a reflected window with symmetric positive weights.
\begin{enumerate}[label=\textup{(\alph*)}]
  \item \(\|v_W\|^2=A_W\), and \(A_W\leq\AZ\).
  \item \(L_Wv_W=0\).  Consequently
    \(\langle v_W,\Xi(L_W)v_W\rangle=A_W\): the invariant probe sees
    nothing of the displacement, and the residual charge is the whole
    window sum.
  \item For every linear map \(L\) on \(\mathbb R^W\),
    \(\langle v_W,(L^+(LL^*)(L^+)^*+\Xi(L))v_W\rangle=A_W\).  Changing the
    probe moves charge between the decoded native currency and the
    residual, but never changes their total.
  \item Let \(P_W\) be the orthogonal projection onto \(\ker L_W\), a probe
    determined by the window alone (it does not depend on the
    displacements).  Then
    \(\langle v_W,\Xi(P_W)v_W\rangle=0\) and
    \(\langle v_W,P_WP_W^*v_W\rangle=A_W\).  In particular, a window
    containing a zero off the critical line has zero residual charge for this probe and strictly positive native currency.
\end{enumerate}
Part (d) shows that residual charge depends on the probe; the
equivalences of \Cref{prop:rh:window-equivalences} concern the
invariant probe \(L_W\).
\end{proposition}

\begin{proof}
(a) is the definition of \(v_W\) and of the extended sum.  (b): the
reflection pairs \(\rho\) with \(1-\rho\), which has the opposite
displacement and the same weight; the only fixed point \(\rho=\frac12\)
has displacement zero.  So the terms of \(L_Wv_W\) cancel, \(v_W\) lies in
\(\ker L_W\), and \(\Xi(L_W)v_W=v_W\).  (c): \(L^+(LL^*)(L^+)^*=L^+L\) is
the orthogonal projection onto \((\ker L)^\perp\), so the operator in
(c) is the identity.  (d): \(\ker P_W=(\ker L_W)^\perp\), which is
orthogonal to \(v_W\) by (b); and \(P_WP_W^*=P_W\) fixes \(v_W\).
\end{proof}

\begin{proposition}[Two elementary limitations]
\label{prop:rh:window-obstructions}
Let \(W\) be a nonempty window.
\begin{enumerate}[label=\textup{(\alph*)}]
  \item \emph{Operator transport bounds for the identity target cannot
    vanish.}  This part does not use the zeros at all, and the same
    argument works for any finite-dimensional target of \(L\).  Let
    \(L:\mathbb R^W\to\mathbb R\) be linear, and suppose
    \(LL^*\preceq\Theta\), \(\Xi(L)\preceq\Omega\) and
    \(L^+\Theta(L^+)^*+\Omega\preceq H\cdot I\) for operators
    \(\Theta,\Omega\) and a real number \(H\).  Then \(H\geq1\).
  \item \emph{The probe cannot be coercive on its own residual.}  If
    \(c\|x\|\leq|L_W x|\) for some \(c>0\) and all \(x\) in the range of
    \(\Xi(L_W)\), then \(L_W\) is injective.  Hence no such \(c\) exists when \(W\) has at
    least two elements.
\end{enumerate}
\end{proposition}

\begin{proof}
(a) Conjugating \(LL^*\preceq\Theta\) by \(L^+\) and adding
\(\Xi(L)\preceq\Omega\) gives
\(I=L^+L+\Xi(L)\preceq L^+\Theta(L^+)^*+\Omega\preceq H\cdot I\), and
\(\mathbb R^W\neq0\).  (b) The range of \(\Xi(L_W)\) is \(\ker L_W\), on
which \(L_W\) vanishes; the inequality forces \(\ker L_W=0\).  A nonzero
linear functional on \(\mathbb R^W\) is injective only if
\(\dim\mathbb R^W=1\).
\end{proof}

Here the \emph{identity target} is the map \(I\) on \(\mathbb R^W\) whose
values the payment is meant to control, and a \emph{transport bound} is
the scalar \(H\) in an inequality of the form in (a), which dominates the
decoded native budget plus the residual bound by \(H\cdot I\).  Part (b) says that coercivity on the residual forces \(L_W\) to be
injective, which fails as soon as \(W\) has at least two elements,
because then \(\ker L_W\neq0\).  Part (a) says only that
the identity operator is not small: a transport
bound of the form in (a), which dominates the identity target as an
operator rather than on the actual displacement vector \(v_W\), cannot
tend to zero.

\begin{proposition}[Payments that restate the Riemann hypothesis]
\label{prop:rh:window-equivalences}
Let \((W_n)\) be an exhaustive sequence of reflected windows with
symmetric positive weights, and write \(v_n=v_{W_n}\), \(L_n=L_{W_n}\),
\(A_n=A_{W_n}\).  Each of the following is equivalent to the Riemann
hypothesis:
\begin{enumerate}[label=\textup{(\alph*)}]
  \item \(\langle v_n,\Xi(L_n)v_n\rangle=0\) for every \(n\);
  \item there are \(b_n\to0\) with \(\langle v_n,\Xi(L_n)v_n\rangle\leq b_n\);
  \item there are operators \(\Omega_n\succeq\Xi(L_n)\) and \(b_n\to0\)
    with \(\langle v_n,\Omega_nv_n\rangle\leq b_n\);
  \item for a fixed real number \(q<1\),
    \(\langle v_n,\Xi(L_n)v_n\rangle\leq q\,A_n\) for every \(n\);
  \item \(\sum_n A_n<\infty\).
\end{enumerate}
\end{proposition}

\begin{proof}
By \Cref{prop:rh:window-identities}(b), the residual charge in (a),
(b) and (d) equals \(A_n\), and in (c) it is at most
\(\langle v_n,\Omega_nv_n\rangle\).  If some zero \(\rho_0\) is off the
critical line, it lies in \(W_n\) for all large \(n\), and then
\(A_n\geq(\Real\rho_0-\frac12)^2>0\) for all large \(n\); this
contradicts each of (a)--(e) (for (d), \(A_n\leq qA_n\) forces
\(A_n=0\); for (e), the terms of a convergent series tend to zero).
Conversely, under the Riemann hypothesis every \(A_n\) is zero, and
(a)--(e) hold with \(\Omega_n=\Xi(L_n)\) and \(b_n=0\).
\end{proof}

The proposition does not say that the Riemann hypothesis is hard; it
says that, on actual zero windows, bounding the residual charge,
bounding any operator above it, asking for a uniform residual gap, or
asking for summable window sums is each the Riemann hypothesis in
another form.  An argument along these lines becomes a proof only if
the bound is derived from zeta-function data by a route that does not
already presuppose it.
 \section{Scope and non-claims}
\label{sec:scope_and_nonclaims}

\paragraph{No unconditional claim.}
No statement about the location of the zeros of \(\zeta\) is proved.
The conditional theorems assume \(\GamSDTCSelberg\), and
\Cref{thm:rh:full-classical} also assumes a classical-zero
identification.  By \Cref{prop:rh:gamma-calibration} and
\Cref{thm:rh:concrete-window-dc}, on the bare shell of all nontrivial zeros with their real parts, and
for windows that eventually
contain every zero, these hypotheses, as types of data, are equivalent
to the Riemann hypothesis; the conditional theorems
are reformulations and provide no evidence for it.

\paragraph{The recognition source is not constructed.}
The intended provenance of \(\GamSDTCSelberg\)---a trace formula of
Selberg type for the zeros of \(\zeta\), realizing the structural law
of the companion paper---is not established, and no trace identity
between spectral and arithmetic data is used.  The shell's name
records that intention only (\Cref{sec:sat_sel_shell}).

\paragraph{Admissibility is not established.}
The framework's formation rules \citep{TsiokosFoundationsII2026} would
require an intended trace-formula realization of the shell to be
admissible; this is not verified, and the shell's audit marker is a
label, not a proof.

\paragraph{The classification is bookkeeping.}
The register of \Cref{sec:rh-aor-instance} is a finite list of labels.
Its completeness (\Cref{thm:rh:aor-instance}) does not imply the
critical-line statement, and the stronger, proof-carrying notion of
realisability of \citet{TsiokosAOR2026} is not established for this
argument.

\paragraph{Other \(L\)-functions and finer questions.}
Only the Riemann zeta function is treated.  Nothing is claimed about
the generalized Riemann hypothesis for Dirichlet, modular or other
\(L\)-functions of the Selberg class
\citep{Selberg1992,KaczorowskiPerelli1999}, about simplicity of the
zeros \citep{Montgomery1973,ConreyGhoshGonek1998}, or about zero
density \citep{Montgomery1971,Ivic1985}.
 \section{Discussion}
\label{sec:discussion}

\paragraph{What the route requires.}
The duality-confinement route reduces the Riemann hypothesis to upper
bounds tending to zero for the window sums \(A_W\), or for operators
dominating the associated ledgers.  For windows closed under the
reflection \(s\mapsto1-\bar s\) (\Cref{thm:rh:concrete-window-dc}), or under the map \(s\mapsto1-s\) with invariant weights
(\Cref{prop:rh:window-equivalences}), and along windows that eventually
contain every zero, the existence of such bounds is the Riemann
hypothesis itself, whatever form the bounds take among those
considered.  What would be new is a source of bounds: an argument, for
instance from an explicit formula or a trace formula, that produces
them from zeta-function data without presupposing the conclusion.
\Cref{prop:rh:window-obstructions} records the elementary fact that an operator-level transport bound for
the identity target, of the form stated there, cannot tend to zero on a
nonempty window, so a payment of that kind must be evaluated on the
actual displacement vector.

\paragraph{Relation to classical programs.}
Several classical approaches seek a positivity statement equivalent to
the Riemann hypothesis: Weil's criterion, which asks for positivity of
a quadratic functional built from the explicit formula
\citep{Weil1952,Bombieri2000}, and Li's criterion \citep{Li1997}; the Nyman--Beurling criterion
\citep{Nyman1950,Beurling1955,BaezDuarte2003}; de Branges' spaces of
entire functions \citep{deBranges1968,deBranges1986}; and spectral
interpretations of the zeros \citep{BerryKeating1990,Connes1999}.  The window sums of this paper are also positivity data, but of a much
cruder kind: they are the squared displacements themselves, and, unlike Weil's functional, which is built from the explicit
formula relating zeros to primes, they involve no primes.  A bound
for them that did not presuppose the Riemann hypothesis would presumably
have to come from a deeper source of this kind; nothing in the present paper
supplies one.

\paragraph{Relation to the companion papers.}
The duality-confinement theorem is proved in \citet{TsiokosSDTC2026};
this paper is its application to the zeros of \(\zeta\), and
\Cref{thm:rh:concrete-window-dc} is the zeta-specific case of the
converse proved there.  The residual operator of \Cref{sec:finite_windows} is the
identity-audit case of the adequacy residual of
\citet{TsiokosXiCompanion2026}, which in the companion Navier--Stokes
paper \citep{TsiokosNavier2026} is applied to a physical target; the
metatheory paper \citep{TsiokosMetaTheory2026} relates the two
applications through a common Gaussian observation.  None of these
connections supplies a payment for either problem.

\paragraph{Further directions.}
Three questions remain open: to construct a recognition source from a
trace formula, which would turn \Cref{thm:rh:full-classical} into a
proof; to verify the framework's admissibility conditions for such a
construction; and to treat other \(L\)-functions, each of which needs
its own ledger and source.
 \section{Conclusion}
\label{sec:conclusion}

The Riemann hypothesis is the vanishing of the positive ledger
\(\AZ=\sum_\rho m_\rho(\Real\rho-\frac12)^2\) of squared displacements of
the nontrivial zeros from the fixed line of the involution
\(s\mapsto1-\bar s\).  If a recognition source supplies domination
records of vanishing trace together with a bridge encoding these displacements in an auxiliary
readout, the duality-confinement theorem places every represented zero on the
critical line of its coordinates, and, with a map from coordinates to the reals that sends \(\frac12\) to
\(\frac12\) and carries the represented zeros onto the real parts of all
nontrivial zeros, the Riemann hypothesis follows.  On any shell this
hypothesis is equivalent to the critical-line statement for the
represented zeros, hence on the bare shell to the Riemann hypothesis, and so
is its concrete form on finite windows closed under the reflection
\(s\mapsto1-\bar s\), with unit weights, that eventually contain every
zero.  For windows closed under \(s\mapsto1-s\), with invariant weights, an
operator-level transport bound for the full identity target cannot
tend to zero on a nonempty window, and several natural forms of a
payment for the window sums along windows that eventually contain
every zero are again equivalent to the Riemann hypothesis.  These facts are elementary; they use analytic properties of \(\zeta\)
such as the functional equation and the zero-free half-plane
\(\Real s\geq1\), but no quantitative or prime-side input.  The paper thereby
describes what a proof along this route would need, shows that the
obvious candidates are no easier than the Riemann hypothesis, and does
not claim to supply one.  The mathematical statements are formalized in Lean~4 with Mathlib
(\Cref{app:formalization}).

\appendix
\crefalias{section}{appendix}
\section{Formal verification}
\label{app:formalization}

The theorems, propositions and definitions of this paper are formally
verified in the Lean~4 proof assistant (version 4.28.0) with the Mathlib library (revision
\leanid{8f9d9cff6bd728b17a24e163c9402775d9e6a365}).  The Riemann
hypothesis in the conclusions is Mathlib's predicate
\texttt{RiemannHypothesis}: every zero of \texttt{riemannZeta} that is
not a trivial zero and not equal to \(1\) has real part \(\frac12\).
This appendix describes how the objects of the paper are represented
and lists the Lean theorem behind each numbered result.

\subsection*{Representation}

The actual nontrivial zeros are the complex numbers \(s\) with
\texttt{completedRiemannZeta}\,\(s=0\) and \(0<\Real s<1\); the
anti-invariant ledger is a sum in \([0,\infty]\) over them with arbitrary
positive natural weights.  A shell is a Lean structure whose zero ledger
uses an abstract real coordinate; its trace, quotient, visibility and
audit fields are arbitrary types without laws, as stated in
\Cref{sec:sat_sel_shell}.  The bare shell instantiates the coordinate as
\(\mathbb R\) and the zero ledger as the actual nontrivial zeros with
unit weights.  The recognition source \(\GamSDTCSelberg\) is a Lean
structure whose fields are the data of
\Cref{obl:rh:gamma-sdtc-selberg}; it is not an axiom, and every theorem
that uses it takes an element of it (data witnessing the hypothesis)
as an explicit argument.  Its
auxiliary ledger is expressed in the abstract, library-free interface
of the first formulation of the companion paper, which also carries some unused structure (weights, a response involution, a
quantitative separation modulus).  The formal and the written versions
of the hypothesis are equivalent: each holds for a shell exactly when
every represented zero has coordinate \(\frac12\), by
\Cref{prop:rh:gamma-calibration} and its formal counterpart.  The concrete window form
(\Cref{thm:rh:concrete-window-dc}) uses instead the Mathlib operator
formulation of the companion paper's Version~3, with positive operators
on \(\mathbb R\) and the linear-algebra trace.

\subsection*{What is not formally verified}

The classical facts cited in the text about the functional equation,
the location of zeros of \(\zeta\) outside the critical strip and the
countability of the zeros are taken from, or derived from, Mathlib's
formal development of the zeta function wherever the proofs use them.  Not verified are: \Cref{rem:rh:aor-partial-status}, which is a scope
statement about the stronger notion of realisability; the framework
material of \Cref{sec:framework}; and the discussion of
provenance and admissibility.  The identification, stated in
\Cref{sec:finite_windows}, of the companion paper's residual and
decoder at identity audit with \(\Xi(L)\) and \(L^+\) is formally
verified (see \Cref{prop:rh:window-identities} in the table).

\Needspace{16\baselineskip}
\subsection*{Correspondence}

{\footnotesize
\begin{longtable}{@{}>{\raggedright\arraybackslash}p{0.21\linewidth}
  >{\raggedright\arraybackslash}p{0.74\linewidth}@{}}
\caption{Numbered results and the Lean declarations that prove or
realize them; names are relative to the namespace
\texttt{SixBirdsDualityConfinement.RH}.}\label{tab:rh-formalization-coverage}\\
\toprule
Result & Lean declarations \\
\midrule
\endfirsthead
\toprule
Result & Lean declarations (continued) \\
\midrule
\endhead
\bottomrule
\endfoot
\Cref{def:rh:fe-involution,def:rh:psi-minus-rh}
  & \leanid{ActualInvolution.J}, \leanid{ActualInvolution.J_involutive},
    \leanid{ActualInvolution.J_fixed_iff}, \leanid{ActualInvolution.J_zero};
    shell form \leanid{Involution.feInvolution}, \leanid{Involution.psiMinusRh} \\ \addlinespace
\Cref{def:rh:nontrivial-zero-ledger,def:rh:anti-invariant-zero-ledger}
  & \leanid{ClassicalZeroLedger.NontrivialZero},
    \leanid{ClassicalZeroLedger.antiInvariantSum};
    shell form \leanid{ZeroLedger.NontrivialZeroLedger},
    \leanid{AntiInvariantZeroLedger.AntiInvariantZeroLedger} \\ \addlinespace
\Cref{def:rh:sat-sel-shell}
  & \leanid{SatSelShell.SatSelShell}; bare shell
    \leanid{ClassicalBareShell.bareShell} \\ \addlinespace
\Cref{thm:rh:translation-T,thm:rh:translation-T-forward,thm:rh:translation-T-reverse}
  & \leanid{ClassicalZeroLedger.antiInvariantSum_eq_zero_iff},
    \leanid{ClassicalZeroLedger.antiInvariantSum_eq_zero_iff_criticalStripRH},
    \leanid{FullClassicalRHBridge.criticalStripRH_iff_riemannHypothesis};
    shell form \leanid{TranslationT.translationT},
    \leanid{TranslationT.translationTForward},
    \leanid{TranslationT.translationTReverse} \\ \addlinespace
\Cref{obl:rh:gamma-sdtc-selberg}
  & \leanid{RHConditional.GammaSdtcSelberg} (a structure; its elements are the data (i)--(iv)) \\ \addlinespace
\Cref{prop:rh:gamma-calibration}
  & \leanid{BareSourceConverse.gamma_nonempty_iff_shell_criticalLine},
    \leanid{BareSourceConverse.gamma_nonempty_iff_criticalStripRH},
    \leanid{BareSourceConverse.gammaOfRH} \\ \addlinespace
\Cref{thm:rh:dc-master-applied}
  & \leanid{DCMasterApplied.dcMasterAppliedRealTrace} \\ \addlinespace
\Cref{thm:rh:conditional}
  & \leanid{RHConditional.rhConditional},
    \leanid{RHConditional.rhConditionalClassical} \\ \addlinespace
\Cref{thm:rh:full-classical}
  & \leanid{FullClassicalRHBridge.rhConditionalFullClassical},
    \leanid{RHConditional.ClassicalZeroIdentification} \\ \addlinespace
\Cref{thm:rh:concrete-window-dc}
  & \leanid{ConcreteZeroWindowDC.windowRecords_family_iff_riemannHypothesis}, \leanid{ConcreteZeroWindowDC.movingWindowRecords_iff_riemannHypothesis};
    \leanid{ConcreteZeroWindowDC.windowPsi_separates},
    \leanid{ConcreteZeroWindowDC.weightedLedger_eq_energy},
    \leanid{ConcreteZeroWindowDC.windowRecords_iff_fixed} \\ \addlinespace
\Cref{def:rh:aor-sel-instance} and
\Cref{thm:rh:aor-mechanical-records,thm:rh:aor-recognition-discharge,thm:rh:aor-gamma-bridge-discharge,thm:rh:aor-auditL-discharge,thm:rh:aor-dc-master-import-discharge,thm:rh:aor-real-coordinate-discharge,thm:rh:aor-translation-interface-discharge}
  & \leanid{AORInstance.assembledCarrier}, \leanid{AORInstance.aorMechanicalRecords},
    \leanid{AORInstance.aorRecognitionDischarge},
    \leanid{AORInstance.aorGammaBridgeDischarge},
    \leanid{AORInstance.aorAuditLDischarge},
    \leanid{AORInstance.aorDCMasterImportDischarge},
    \leanid{AORInstance.aorRealCoordinateDischarge},
    \leanid{AORInstance.aorTranslationInterfaceDischarge} (literal lists) \\ \addlinespace
\Cref{thm:rh:aor-instance}
  & \leanid{AORInstance.aorInstance} \\ \addlinespace
\Cref{prop:rh:window-identities}
  & \leanid{FiniteZeroWindows.displacementVector_norm_sq},
    \leanid{FiniteZeroWindows.finite_window_le_global_classical_ledger},
    \leanid{FiniteZeroWindows.invariantProbe_displacement_zero},
    \leanid{FiniteZeroWindows.fixedTarget_window_quadratic_eq_windowEnergy},
    \leanid{FiniteZeroWindows.window_complete_currency_eq_windowEnergy},
    \leanid{FiniteZeroWindows.blindSubspaceNative_actual_residual_zero},
    \leanid{FiniteZeroWindows.blindSubspaceNative_actual_currency_eq_windowEnergy},
    \leanid{FiniteZeroWindows.blindSubspaceNative_offLine_control};
    residual as projection and decoder as pseudoinverse:
    \leanid{SixBirdsNeedles.XiCore.auditResidual_id_eq_hilbertKernelXi},
    \leanid{SixBirdsNeedles.XiCore.adjoint_pinv_gram_eq_pinv} (from the
    companion Needles library), \leanid{ConcreteWindowXI.identity_auditDecoder_comp} \\ \addlinespace
\Cref{prop:rh:window-obstructions}
  & \leanid{FiniteZeroWindows.window_fullReadout_transport_floor},
    \leanid{FiniteZeroWindows.native_coercive_on_own_residual_forces_injective},
    \leanid{ConcreteWindowXI.invariantProbe_not_injective}, \leanid{ConcreteWindowXI.invariantProbe_no_residual_coercivity} \\ \addlinespace
\Cref{prop:rh:window-equivalences}
  & \leanid{FiniteZeroWindows.exhaustive_fixedTarget_xi_zero_iff_criticalStripRH},
    \leanid{FiniteZeroWindows.exists_vanishing_fixedTarget_xi_budget_iff_criticalStripRH},
    \leanid{FiniteZeroWindows.exists_operatorResidualBudget_iff_criticalStripRH},
    \leanid{FiniteZeroWindows.subunit_fixedTarget_xi_gap_iff_criticalStripRH},
    \leanid{FiniteZeroWindows.summable_windowEnergy_iff_criticalStripRH};
    with \leanid{FullClassicalRHBridge.criticalStripRH_iff_riemannHypothesis} \\
\end{longtable}
}

In \Cref{sec:finite_windows} the formal statements use the reflection
\(\rho\mapsto1-\rho\) and positive weights symmetric under it; those for
\Cref{thm:rh:concrete-window-dc} use the involution \(\JL\) and unit
weights.  Statements whose Lean conclusion is the critical-strip form of
the Riemann hypothesis are combined with
\leanid{FullClassicalRHBridge.criticalStripRH_iff_riemannHypothesis}.

\subsection*{Trust base and availability}

Every declaration listed depends only on the standard axioms
\texttt{propext}, \texttt{Classical.choice} and \texttt{Quot.sound}; no
\texttt{sorry} and no additional axioms occur.  The development is
available at
\url{https://github.com/ioannist/six-birds-duality-confinement}; it
also contains further finite-window and Gaussian-observation results
not used in this paper.
 \section{Version history}
\label{app:version-history}
\begingroup\small

Every numbered result of Version~1 keeps its number and place in later
versions, and its name except where noted below.

\subsection*{Version 3 (October 1, 2026)}

\begin{itemize}
  \item The paper was rewritten for readers outside the framework; it
    now says plainly that the conditional theorem is a reformulation of
    the Riemann hypothesis, not evidence for it.  Discussion of the formal
    development was moved to \Cref{app:formalization}; the numbered
    lists of non-claims and ``no-smuggling gates'' and the reference to
    a ``V-Differential'' classification, which is not defined in the
    cited literature, were removed; the account of a three-route
    derivation attempt for the recognition source was removed.
  \item \Cref{def:rh:nontrivial-zero-ledger} now allows arbitrary
    positive weights, with multiplicities as the natural choice;
    Version~2 required multiplicities, which the formal development does
    not use.  \Cref{thm:rh:translation-T} is stated for the actual
    nontrivial zeros and includes the equivalence with the Riemann
    hypothesis; the shell version is described in the text.
  \item \Cref{obl:rh:gamma-sdtc-selberg} is now stated in mathematical
    terms and displayed as a hypothesis (it was labelled
    ``Obligation'').  The new \Cref{prop:rh:gamma-calibration} states that for every shell it
    is equivalent to the critical-line statement for the represented
    zeros, hence on the bare shell to the Riemann hypothesis; Version~2
    stated the bare-shell case in prose.  \Cref{def:rh:sat-sel-shell}
    now states the coordinate system and zero ledger of a shell
    explicitly.
  \item The new \Cref{thm:rh:concrete-window-dc} applies the
    duality-confinement theorem, in the concrete operator form of the
    companion paper's Version~3, to finite reflection-closed windows of
    actual zeros that eventually contain every zero, and proves that the
    resulting domination records exist if and only if the Riemann
    hypothesis holds.
  \item The AOR classification (\Cref{sec:rh-aor-instance}) was
    condensed.  Items~8.2--8.8, which record classification labels and
    were displayed as theorems, are now displayed as ``Records'' and
    describe the recorded lists exactly; \Cref{thm:rh:aor-instance} now
    includes the non-emptiness of the non-claim list.  The section was
    renamed ``A classification of the hypotheses''.
  \item The finite-window material of the Version~2 discussion was
    reorganized into the new \Cref{sec:finite_windows}, with the new
    \Cref{prop:rh:window-identities,prop:rh:window-obstructions,prop:rh:window-equivalences}
    and complete proofs.
  \item Names shortened: \Cref{def:rh:sat-sel-shell} (formerly
    ``Saturated completed Selberg trace shell interface''),
    \Cref{thm:rh:dc-master-applied} (formerly ``Duality-confinement
    master theorem, applied to'' the shell),
    \Cref{thm:rh:full-classical} (formerly ``Full classical conditional
    RH''), Records~8.2--8.8 (formerly ``Mechanical AOR records'',
    ``Recognition-source discharge'', ``Bridge-field discharge'',
    ``Admissibility-audit status record'', ``Cross-paper master-theorem
    discharge'', ``Typed real-coordinate discharge'',
    ``Translation-interface discharge''), \Cref{thm:rh:aor-instance} (formerly ``RH local-register theorem''),
    and \Cref{rem:rh:aor-partial-status} (formerly ``Partial status of
    the AOR-instance wrapper'').
  \item Cross-references now name the kind of result they point to;
    bibliography entries were corrected.
\end{itemize}

\subsection*{Version 2 (September 28, 2026)}

Version~1 described the argument as a refinement-stable
audited-operational-realisability instance and identified the shell's
real-part coordinate with the classical real part by construction.
Version~2 replaced the first claim by the local-register statement,
which is what was proved; made the identification an explicit
hypothesis, constructed it for the bare shell, and proved the passage
from the critical strip to the full Riemann hypothesis; recorded that on
the bare shell the recognition source exists exactly under the Riemann
hypothesis; added the finite-window analysis and the comparison with
the companion Navier--Stokes and metatheory papers; and moved the formal
development from a library-free encoding to Mathlib, correcting the
description of what it proves.
\par\endgroup

\clearpage
\bibliographystyle{plainnat}
\begin{thebibliography}{31}
\providecommand{\natexlab}[1]{#1}
\providecommand{\url}[1]{\texttt{#1}}
\expandafter\ifx\csname urlstyle\endcsname\relax
  \providecommand{\doi}[1]{doi: #1}\else
  \providecommand{\doi}{doi: \begingroup \urlstyle{rm}\Url}\fi

\bibitem[B{\'a}ez-Duarte(2003)]{BaezDuarte2003}
Luis B{\'a}ez-Duarte.
\newblock A strengthening of the {Nyman}--{Beurling} criterion for the {Riemann} hypothesis.
\newblock \emph{Rendiconti Lincei. Matematica e Applicazioni}, 14\penalty0 (1):\penalty0 5--11, 2003.

\bibitem[Berry and Keating(1990)]{BerryKeating1990}
M.~V. Berry and J.~P. Keating.
\newblock A rule for quantizing chaos?
\newblock \emph{Journal of Physics A: Mathematical and General}, 23\penalty0 (21):\penalty0 4839--4849, 1990.
\newblock \doi{10.1088/0305-4470/23/21/024}.

\bibitem[Beurling(1955)]{Beurling1955}
Arne Beurling.
\newblock A closure problem related to the {Riemann} zeta-function.
\newblock \emph{Proceedings of the National Academy of Sciences of the United States of America}, 41\penalty0 (5):\penalty0 312--314, 1955.
\newblock \doi{10.1073/pnas.41.5.312}.

\bibitem[Bombieri(2000)]{Bombieri2000}
Enrico Bombieri.
\newblock Remarks on {Weil}'s quadratic functional in the theory of prime numbers, {I}.
\newblock \emph{Rendiconti Lincei. Matematica e Applicazioni}, 11\penalty0 (3):\penalty0 183--233, 2000.

\bibitem[Connes(1999)]{Connes1999}
Alain Connes.
\newblock Trace formula in noncommutative geometry and the zeros of the {Riemann} zeta function.
\newblock \emph{Selecta Mathematica, New Series}, 5\penalty0 (1):\penalty0 29--106, 1999.
\newblock \doi{10.1007/s000290050042}.

\bibitem[Conrey et~al.(1998)Conrey, Ghosh, and Gonek]{ConreyGhoshGonek1998}
J.~Brian Conrey, Amit Ghosh, and Steve~M. Gonek.
\newblock Simple zeros of the {Riemann} zeta-function.
\newblock \emph{Proceedings of the London Mathematical Society}, 76\penalty0 (3):\penalty0 497--522, 1998.
\newblock \doi{10.1112/S0024611598000306}.

\bibitem[de~Branges(1968)]{deBranges1968}
Louis de~Branges.
\newblock \emph{{Hilbert} Spaces of Entire Functions}.
\newblock Prentice-Hall, Englewood Cliffs, NJ, 1968.

\bibitem[de~Branges(1986)]{deBranges1986}
Louis de~Branges.
\newblock The {Riemann} hypothesis for {Hilbert} spaces of entire functions.
\newblock \emph{Bulletin of the American Mathematical Society}, 15\penalty0 (1):\penalty0 1--17, 1986.
\newblock \doi{10.1090/S0273-0979-1986-15429-7}.

\bibitem[Edwards(1974)]{Edwards1974}
Harold~M. Edwards.
\newblock \emph{{Riemann}'s Zeta Function}.
\newblock Academic Press, New York, 1974.

\bibitem[Hejhal(1976)]{Hejhal1976}
Dennis~A. Hejhal.
\newblock \emph{The {Selberg} Trace Formula for {PSL}(2,{R}), Volume 1}, volume 548 of \emph{Lecture Notes in Mathematics}.
\newblock Springer, Berlin, Heidelberg, 1976.
\newblock \doi{10.1007/BFb0079608}.

\bibitem[Ivi{\'c}(1985)]{Ivic1985}
Aleksandar Ivi{\'c}.
\newblock \emph{The {Riemann} Zeta-Function: The Theory of the {Riemann} Zeta-Function with Applications}.
\newblock John Wiley \& Sons, New York, 1985.
\newblock ISBN 978-0-471-80634-9.

\bibitem[Iwaniec(2002)]{Iwaniec2002}
Henryk Iwaniec.
\newblock \emph{Spectral Methods of Automorphic Forms}, volume~53 of \emph{Graduate Studies in Mathematics}.
\newblock American Mathematical Society, Providence, RI, second edition, 2002.
\newblock ISBN 978-0-8218-3160-1.

\bibitem[Kaczorowski and Perelli(1999)]{KaczorowskiPerelli1999}
Jerzy Kaczorowski and Alberto Perelli.
\newblock The {Selberg} class: A survey.
\newblock In K{\'a}lm{\'a}n Gy{\H{o}}ry, Henryk Iwaniec, and Jerzy Urbanowicz, editors, \emph{Number Theory in Progress}, volume~2, pages 953--992. Walter de Gruyter, Berlin and New York, 1999.
\newblock \doi{10.1515/9783110285581.953}.

\bibitem[Li(1997)]{Li1997}
Xian-Jin Li.
\newblock The positivity of a sequence of numbers and the {Riemann} hypothesis.
\newblock \emph{Journal of Number Theory}, 65\penalty0 (2):\penalty0 325--333, 1997.
\newblock \doi{10.1006/jnth.1997.2137}.

\bibitem[Meyer(2005)]{Meyer2005}
Ralf Meyer.
\newblock A spectral interpretation for the zeros of the {Riemann} zeta function.
\newblock In \emph{Mathematisches Institut, Georg-August-Universit{\"a}t G{\"o}ttingen: Seminars Winter Term 2004/2005}, pages 117--137. Universit{\"a}tsdrucke G{\"o}ttingen, G{\"o}ttingen, 2005.

\bibitem[Montgomery(1971)]{Montgomery1971}
Hugh~L. Montgomery.
\newblock \emph{Topics in Multiplicative Number Theory}, volume 227 of \emph{Lecture Notes in Mathematics}.
\newblock Springer, Berlin, Heidelberg, 1971.
\newblock \doi{10.1007/BFb0060851}.

\bibitem[Montgomery(1973)]{Montgomery1973}
Hugh~L. Montgomery.
\newblock The pair correlation of zeros of the zeta function.
\newblock In \emph{Analytic Number Theory}, volume~24 of \emph{Proceedings of Symposia in Pure Mathematics}, pages 181--193, Providence, RI, 1973. American Mathematical Society.

\bibitem[Nyman(1950)]{Nyman1950}
Bertil Nyman.
\newblock \emph{On the One-Dimensional Translation Group and Semi-Group in Certain Function Spaces}.
\newblock PhD thesis, University of Uppsala, 1950.

\bibitem[Riemann(1859)]{Riemann1859}
Bernhard Riemann.
\newblock {Ueber die Anzahl der Primzahlen unter einer gegebenen Gr{\"o}sse}.
\newblock \emph{Monatsberichte der K{\"o}niglich Preu{\ss}ischen Akademie der Wissenschaften zu Berlin}, pages 671--680, 1859.

\bibitem[Selberg(1956)]{Selberg1956}
Atle Selberg.
\newblock Harmonic analysis and discontinuous groups in weakly symmetric {Riemannian} spaces with applications to {Dirichlet} series.
\newblock \emph{The Journal of the Indian Mathematical Society}, 20:\penalty0 47--87, 1956.

\bibitem[Selberg(1992)]{Selberg1992}
Atle Selberg.
\newblock Old and new conjectures and results about a class of {Dirichlet} series.
\newblock In Enrico Bombieri et~al., editors, \emph{Proceedings of the Amalfi Conference on Analytic Number Theory}, pages 367--385, Salerno, 1992. Universit{\`a} di Salerno.

\bibitem[Titchmarsh(1986)]{TitchmarshHeathBrown1986}
E.~C. Titchmarsh.
\newblock \emph{The Theory of the {Riemann} Zeta-Function}.
\newblock Oxford University Press, Oxford, second edition, 1986.
\newblock ISBN 978-0-19-853369-6.
\newblock Revised by D. R. Heath-Brown.

\bibitem[Tsiokos(2026{\natexlab{a}})]{TsiokosAOR2026}
Ioannis Tsiokos.
\newblock Audited operational realisability: Canonical inference closure and evidence-preserving composition.
\newblock \doi{10.5281/zenodo.23086977}, 2026{\natexlab{a}}.
\newblock Version 3, October 1, 2026.

\bibitem[Tsiokos(2026{\natexlab{b}})]{TsiokosFoundationsI2026}
Ioannis Tsiokos.
\newblock {Six Birds: Foundations of Emergence Calculus}.
\newblock \href{https://arxiv.org/abs/2602.00134}{arXiv:2602.00134}, 2026{\natexlab{b}}.

\bibitem[Tsiokos(2026{\natexlab{c}})]{TsiokosFoundationsII2026}
Ioannis Tsiokos.
\newblock {Six Birds Foundations II: Admissibility, Meta-Theory, and the Exact-Six Program}.
\newblock \doi{10.5281/zenodo.19672278}, 2026{\natexlab{c}}.

\bibitem[Tsiokos(2026{\natexlab{d}})]{TsiokosFoundationsIII2026}
Ioannis Tsiokos.
\newblock {Six Birds Foundations III: A Finite Audited Interaction Calculus for SBT}.
\newblock \doi{10.5281/zenodo.20124595}, 2026{\natexlab{d}}.

\bibitem[Tsiokos(2026{\natexlab{e}})]{TsiokosMetaTheory2026}
Ioannis Tsiokos.
\newblock Three conditional {Clay}-problem closures: A shared {Gaussian} observation and a common vocabulary.
\newblock \doi{10.5281/zenodo.23086999}, 2026{\natexlab{e}}.
\newblock Version 4, October 1, 2026.

\bibitem[Tsiokos(2026{\natexlab{f}})]{TsiokosNavier2026}
Ioannis Tsiokos.
\newblock A conditional spectral-closure theorem for three-dimensional {Navier--Stokes}.
\newblock \doi{10.5281/zenodo.23086987}, 2026{\natexlab{f}}.
\newblock Version 4, October 1, 2026.

\bibitem[Tsiokos(2026{\natexlab{g}})]{TsiokosSDTC2026}
Ioannis Tsiokos.
\newblock Self-dual trace confinement: A {Six Birds} structural law for formed closures under involutive self-duality.
\newblock \doi{10.5281/zenodo.23086972}, 2026{\natexlab{g}}.
\newblock Version 3, October 1, 2026.

\bibitem[Tsiokos(2026{\natexlab{h}})]{TsiokosXiCompanion2026}
Ioannis Tsiokos.
\newblock Adequacy residuals and blind-spot currency.
\newblock \doi{10.5281/zenodo.23086980}, 2026{\natexlab{h}}.
\newblock Version 4, October 1, 2026.

\bibitem[Weil(1952)]{Weil1952}
Andr{\'e} Weil.
\newblock Sur les ``formules explicites'' de la th{\'e}orie des nombres premiers.
\newblock \emph{Meddelanden fr{\aa}n Lunds Universitets Matematiska Seminarium}, Tome suppl{\'e}mentaire:\penalty0 252--265, 1952.

\end{thebibliography}


\end{document}
